\section{Clifford Error Analysis Proofs}
\label{sec:clifford_error_analysis_proofs}
In this section we prove \cref{theorem:probability_trivial_syndrome,prop:resulting_fidelity}.
First note the following.
We can write the density operator of
the logical state corresponding to \cref{eq:pauli_decomposition} in terms of
	$\mathcal G$ (defining the code)
	and the coefficients $`{\alpha_L}_L$ (defining the logical state):
\begin{align}
\overline{\rho}
=\upPi_{\mathcal G} 2^{-k} \sum_{L \in \mathcal{P}_k} \alpha_L \overline{L}.
\label{eq:density_logical_psi}
\end{align}
Here, $\upPi_{\mathcal G}$ projects onto the codespace of $\mathcal G$,
and the other term projects onto the specific logical state within the codespace.
Note $\upPi_{\mathcal G}$ commutes with each $\overline{L}$.

\subsection{Proof of \texorpdfstring
{\cref{theorem:probability_trivial_syndrome}}
{Theorem~\ref{theorem:probability_trivial_syndrome}}}
Let
$\mathcal A:=\langle\mathcal G\rangle$,
$\mathcal B:=\langle E^\dagger\mathcal G E\rangle$,
$\mathcal K:=\mathcal A\cap\mathcal B$,
and $r :=n-k$.
The probability to observe a trivial syndrome
is the expectation value of $\upPi_{\mathcal G}$
[recall \cref{eq:pj_defn}]:
\begin{align}
	\pr(\v s =\v 0)
	&=\tr`*(\upPi_{\mathcal G} E\overline{\rho} E^\dag) \notag \\
	&=\tr`*(\upPi_{E^\dag \mathcal G E} \overline{\rho}) \notag \\
	&\overset{\eqref{eq:density_logical_psi}}=
	2^{-k}\sum_{L\in\mathcal P_k}\alpha_L \tr\left(
	\upPi_{E^\dagger \mathcal G E}\upPi_{\mathcal G} \overline L
	\right).
	\label{eq:sum_of_traces_fully_expanded}
\end{align}
If $\mathcal B\cap(-\mathcal A)\ne\varnothing$,
then $\upPi_{E^\dagger \mathcal G E}\upPi_{\mathcal G}$
contains $(\I -P)(\I +P) =0$ for some Pauli $P \in \mathcal A$,
so the probability is zero.
Assume otherwise.

Multiply out the two stabiliser projectors:
\begin{equation}
\label{eq:trace_triple}
\pr(\v s =\v 0)
=2^{-2r -k} \sum_{L \in \mathcal P_k} \alpha_L
\sum_{a \in \mathcal A} \sum_{b \in \mathcal B}
\tr`\big(ba \overline{L}).
\end{equation}
Since nontrivial Paulis are traceless,
we need only consider products
$ba \overline{L}$
that are proportional to $\I$.
Thus, each $L$ term contributes iff \cref{eq:logical_membership} holds;
equivalently,
iff $\mathcal B\cap\omega_L\overline L\mathcal A$ is nonempty
for some sign $\omega_L\in\{\pm1\}$.
(This sign is real:
	$\overline{L}$ and all Paulis in $\mathcal A$ and $\mathcal B$ are Hermitian,
	and $\overline{L}$ commutes with each Pauli in $\mathcal A$.
The sign is also unique:
	if $b=\omega\overline L a$ and
	$b'=-\omega\overline L a'$ with
	$b,b'\in\mathcal B$ and $a,a'\in\mathcal A$,
	then $bb'=-aa'\in\mathcal B\cap(-\mathcal A)$,
	a contradiction.)
If the intersection is nonempty and contains $b_L$,
then it is a coset of $\mathcal K$:
\begin{equation}
\mathcal B\cap\omega_L\overline L\mathcal A=b_L\mathcal K,
\end{equation}
so it has $2^{\rk\mathcal K}$ elements.
Each surviving trace contributes
$\tr \omega_L \I= \omega_L 2^n =\omega_L 2^{r +k}$,
and therefore the total prefactor is
$2^{-r+\rk\mathcal K}$,
where
$r-\rk\mathcal K
=\rk\langle\mathcal G\cup E^\dagger\mathcal G E\rangle-r
=s$. \qed

\subsection{Proof of \texorpdfstring
{\cref{prop:resulting_fidelity}}
{Proposition~\ref{prop:resulting_fidelity}}}
The fidelity of the resulting state $\ket{\Omega}$
to $\ket{\overline{T}}$ is $|`<\overline{T}|\Omega>|^2$.
Using the notation of \cref{eq:pj_defn}:
\begin{align}
	\ket{\Omega}
	&=\frac1{\sqrt{\pr(\v s =\v 0)}}
	\upPi_{\mathcal G}
	E\ket{\overline{T}} \notag \\
	|`<\overline{T}|\Omega>|^2 \pr(\v s =\v 0)
	&=\tr`\big(E^\dag \proj{\overline{T}}E \proj{\overline{T}}).
\end{align}
Similar to \cref{eq:density_logical_psi},
we can write $\proj{\overline{T}} =\upPi_{\mathcal G} \frac12(\I +\overline{H_+})$,
where $\overline{H_+}$ here refers to
the logical representative from \cref{eq:logical_H_is_transversal}.
Let $a=0$ if $E$ commutes with $\overline{H_+}$
and $a=1$ if it anticommutes,
then $E^\dag \overline{H_+}E =(-1)^a \overline{H_+}$.
Substitute these two facts in:
\begin{align}
|`<\overline{T}|\Omega>|^2 \pr(\v s =\v 0)
&=\tr`\bigg[E^\dag \upPi_{\mathcal G} E \frac{\I +(-1)^a \overline{H_+}}2 \upPi_{\mathcal G} \frac{\I +\overline{H_+}}2]
\notag \\
&=
\tr`\bigg[\upPi_{E^\dagger \mathcal G E}\upPi_{\mathcal G}
\frac{\I +(-1)^a \overline{H_+}}2
\frac{\I +\overline{H_+}}2].
\end{align}
If $a =1$,
the last two projectors multiply to zero,
so $|`<\overline{T}|\Omega>|^2 =0$.
If $a =0$,
they are the same projector;
since projectors are idempotent,
we keep only one of them in our expression.
That remaining logical projector is multiplied by $\upPi_{\mathcal G}$
i.e.\ acts only on the codespace,
so we can replace it with
the logically equivalent projector
$\sum_{L \in \mathcal P_1} \alpha_L \overline{L}$
from \cref{eq:density_logical_psi}
(here $\alpha_\I =1$ and $\alpha_X =\alpha_Y =\frac1{\sqrt{2}}$).
Then:
\begin{equation}
|`<\overline{T}|\Omega>|^2 \pr(\v s =\v 0)
=2^{-1}\sum_{L\in\mathcal P_k}\alpha_L \tr\left(
\upPi_{E^\dagger \mathcal G E}\upPi_{\mathcal G} \overline L
\right),
\label{eq:sum_of_traces_T_fully_expanded}
\end{equation}
This is identical to the right-hand side of
\cref{eq:sum_of_traces_fully_expanded},
so $|`<\overline{T}|\Omega>|^2 =1$.

The stated time complexity comes from determining $a$;
this requires checking the commutation of each pair of tensor factors,
which takes $\mathcal O(n)$ time. \qed